\subsection{معرفی و انگیزه}
این مکانیزم بر مبنای مبانی نظری پیشرفته در ادبیات توجه و فشرده‌سازی \cite{ref_ttcompressor_2024} توسعه یافته است.

تجزیه قطار تانسوری (\lr{Tensor-Train Decomposition - TT}) یکی از تکنیک‌های پیشرفته در جبر تانسوری چندبعدی است که برای فشرده‌سازی تانسورهای مرتبه بالا و کاهش پیچیدگی نمایی ابعاد به کار می‌رود. ایده بنیادی این روش، نگاشت یک ماتریس دوبعدی بزرگ به فضایی چندبعدی و سپس نمایش آن به شکل ضرب زنجیره‌ای از هسته‌های تانسوری کم‌مرتبه سه‌بعدی و چهاربعدی است.

موتور فشرده‌سازی \lr{TT Compressor} با هدف جایگزینی لایه‌های متراکم خطی \lr{nn.Linear} و \lr{Conv1D} با ماژول‌های عامل‌بندی‌شده تانسوری به نام \lr{TTLinearLayer} طراحی گردید. با این وجود، تحلیل‌های عمیق تئوری و تجربی در این پروژه نشان داد که اعمال تجزیه قطار تانسوری بر روی ماتریس‌های تبدیل خطی دوبعدی در مدل‌های زبانی ترنسفورمر دارای چالش‌های بنیادین است و به دلیل پدیده درهم‌آمیختگی ابعاد و افت کند مقادیر مقیاس تانسوری نسبت به \lr{SVD} دوبعدی، این ماژول به عنوان یک رویکرد منسوخ‌شده (\lr{Deprecated}) رده‌بندی شده است.
